\chapter{直方图：原子操作与私有化}

\begin{solutionexercise}{1}
\problemheading 假设 DRAM 系统中每次原子操作的总延迟为 100 ns。针对同一个全局内存变量
执行原子操作时，能够获得的最大吞吐量是多少？

\approachheading 同一地址的读--改--写必须串行化，因此理想吞吐量是单次总延迟的倒数。

\answerheading
\[
R=\frac{1}{100\times10^{-9}\ \mathrm{s}}
 =10^7\ \mathrm{atomic\ ops/s}
 =10\ \mathrm{Mops/s}.
\]

\conclusionheading 最大约为每秒 1000 万次原子操作。
\discussionheading
\discussionpoint{从倒数公式走向排队模型。}把同一地址想成只有一个服务窗口，读、修改、写回构成不可重叠的状态转移；在题设的固定服务时间模型里，服务率就是 $\mu=1/L=10^7$ 次/秒。只有当到达过程带有随机性并且长期到达率 $\lambda$ 接近 $\mu$ 时，排队论才进一步预言等待时间急剧增加；若请求严格等间隔，不能无条件宣称存在同样的排队增长。若暂用指数服务时间的 M/M/1 近似，$1/(1-\rho)$ 是平均逗留时间相对服务时间 $L$ 的放大倍数，而纯排队等待为 $W_q=\rho L/(1-\rho)$，两者不能混称。题设把服务时间固定为 $L$，在泊松到达假设下更接近 M/D/1，此时 $W_q=\rho L/[2(1-\rho)]$；真实 GPU 请求的成批到达和调度相关性还会偏离这两个模型。这个区分说明吞吐上界与响应时间不是同一概念，也解释了为何热点计数器常在平均吞吐尚可时先出现长尾延迟。

\discussionpoint{地址级并行如何打破单热点上界。}十兆次每秒只约束同一状态变量，因为不同地址没有“必须看到上一次结果”的共同依赖链。若更新均匀落到 $H$ 个真正独立且由足够多内存分区服务的地址，粗略上界可以增长到 $H/L$；但缓存行映射、分区冲突和有限的原子执行单元会使增长早于理想值饱和。算法因此常先在寄存器、warp、线程块和全局分片中逐级聚合，再提交少量原子更新，这与树形归约、数据库分区聚合以及分布式计数器的分片思想属于同一种“先拆热点、后合并”的设计。例如八个分片的理想值是 80 Mops/s，若实测仅 45 Mops/s，就应继续检查这些地址是否落到少数内存分区。

\discussionpoint{原子性不等于完整的并发契约。}一次原子读--改--写保证该位置不会出现撕裂更新，却不自动把周围普通内存访问组成事务，也不必然提供应用所需的先后可见性；CUDA 的作用域和内存序仍需按生产者--消费者关系选择。即使竞争已经消除，32 位整数计数器仍可能回绕，浮点加法也会因线程到达次序改变舍入路径，例如 $(10^{20}+(-10^{20}))+1$ 与 $10^{20}+((-10^{20})+1)$ 的结果不同。由此可见，“没有数据竞争”“数学结果正确”和“逐位可复现”是三层不同要求，金融累计、科学复现实验或确定性训练必须分别设计验证策略。

\discussionpoint{怎样把微基准变成可解释证据。}只测一个线程规模会把原子管线、普通带宽和调度不足混在一起，无法判断十兆次上界是否真正生效。可信实验应固定更新总数，再系统改变热点地址数、地址到缓存行的分布、输入偏斜和并发块数，并把吞吐曲线与原子请求、缓存扇区和 warp 停顿关联起来；若增加独立地址后近似线性提升，才支持“同址序列化”判断。这样的实验方法也适用于锁、队列尾指针和图遍历 frontier 等共享状态：先构造可控竞争，再用不变量核对结果，最后才解释性能计数器。每组运行还应检查最终计数等于请求总数，否则一个丢更新的“更快”内核会制造完全错误的性能结论。
\variantheading 若总延迟降为 50 ns，同址上限为 $2\times10^7$ ops/s，即 20 Mops/s。
\practiceheading 用 Nsight Compute 的 atomic throughput 与 L2 sectors 在目标 GPU 上验证。
\sourcecheck{https://docs.nvidia.com/cuda/cuda-c-programming-guide/}{NVIDIA CUDA C++ Programming Guide 对原子读--改--写和同址序列化的定义；量纲独立复核，置信度高}
\end{solutionexercise}

\begin{solutionexercise}{2}
\problemheading 某处理器支持在 L2 缓存中执行原子操作。假设每次原子操作在 L2 缓存中需要
4 ns，在 DRAM 中需要 100 ns，并且 90\% 的原子操作命中 L2 缓存。针对同一个
全局内存变量执行原子操作时，近似吞吐量是多少？

\approachheading 对同一地址的串行操作，长期平均服务时间是两种延迟按命中概率的
加权平均，再取倒数；不能把两条独立吞吐率直接相加。

\answerheading
\[
\bar L=0.9(4)+0.1(100)=13.6\ \mathrm{ns},
\qquad
R\approx\frac1{13.6\times10^{-9}}
=7.35294\times10^7\ \mathrm{ops/s}.
\]
即约 $73.5$ Mops/s。

\conclusionheading 近似吞吐量为 $7.35\times10^7$ 次/秒（约 73.5 Mops/s）。
\discussionheading
\discussionpoint{期望延迟为何不要求独立命中。}设第 $i$ 次原子的服务时间 $L_i$ 只可能为 4 ns 或 100 ns，只要长期命中比例是 $0.9$，总时间除以请求数就趋近 $0.9\times4+0.1\times100=13.6$ ns。这里使用的是期望的线性性，并不要求相邻命中事件独立；即使九十次命中集中在一起、十次未命中也集中在一起，固定服务时间的串行模型仍给出相同平均值。真正错误的是平均两条吞吐率，因为那等价于把两类请求交给并行服务台，而题目明确要求同一地址按一个序列完成。相关性可以改变短时间窗口中的方差和尾延迟，却不会在这个固定成本模型中改变长期总服务时间，这两种统计量应分别讨论。

\discussionpoint{什么时候命中顺序才会改变现实吞吐。}真实缓存中的 4 ns 和 100 ns 很少是永远不变的常数：连续未命中可能填满队列，写回会占用链路，其他流量也会改变缓存行驻留和内存分区等待。只有加入这些状态相关服务时间、有限队列或可重叠请求后，命中成簇与否才会影响总耗时，而不仅是影响单次延迟的排列。例如两次 DRAM 请求若能并行在不同分区飞行，交错序列可能比严格同址模型更快；若都争用同一分区，则可能更慢。因而题目答案是明确抽象下的精确期望，工程外推必须先说明新增机制，不能把“相关性”本身误当成加权平均失效的原因。

\discussionpoint{缓存命中与拆分热点是两个优化维度。}让原子停留在 L2 可以缩短每次状态转移，却没有消除同一变量的顺序依赖，所以命中率从低到高最终仍会碰到单地址服务率。若把一个计数器按 key、warp 或线程块拆成多个副本，既可能提高这些缓存行的驻留概率，又创造真正的地址级并行；代价是初始化副本、占用容量并在末尾再做一次合并。这个权衡与数据库分区计数、分布式 sharded counter 十分相似：缓存优化降低一次访问的成本，分片优化改变依赖图，两者必须在模型中分开记账。若 32 个副本最终仍同时提交到一个地址，合并阶段会重新形成短时热点，所以还要把阶段峰值纳入容量设计。

\discussionpoint{由两点估计走向可校准性能模型。}可以先用常驻单热点和主动驱逐两组实验估计快、慢路径，再在不同命中比例下检查 $1/E[L]$ 是否预测实测饱和吞吐。若误差随并发请求数增加而扩大，通常说明服务可以重叠或队列开始饱和；若误差主要随工作集扩大而变化，则应调查容量失效、缓存行迁移和分区映射。进一步可把模型写成测得的条件延迟 $E[L\mid\text{state}]$ 与状态占比之和，并用独立数据集验证，而不是不断为同一条曲线调参数；这种校准--验证流程也是分析缓存、TLB 和远端内存系统的通用方法。
\variantheading 命中率为 $80\%$ 时平均延迟为 23.2 ns，吞吐约 43.1 Mops/s。
\practiceheading 用 Nsight Compute 的 L2 hit rate 和 atomic 指标代入实测值，并按热点分组分析。
\sourcecheck{https://docs.nvidia.com/cuda/cuda-c-programming-guide/}{NVIDIA CUDA C++ Programming Guide 的原子操作语义与缓存层次说明；以平均服务时间而非平均速率作对抗复核}
\end{solutionexercise}

\begin{solutionexercise}{3}
\problemheading 在第 1 题的条件下，假设内核每执行一次原子操作就执行 5 次浮点运算。受
原子操作吞吐量限制时，内核执行能达到的最大浮点吞吐量是多少？

\approachheading 每秒原子操作数乘每次原子操作关联的浮点运算数。

\answerheading
\[
(10^7\ \mathrm{atomic\ ops/s})(5\ \mathrm{FLOP/atomic\ op})
=5\times10^7\ \mathrm{FLOP/s}
=0.05\ \mathrm{GFLOP/s}.
\]

\conclusionheading 最大浮点吞吐量为 50 MFLOP/s（0.05 GFLOP/s）。
\discussionheading
\discussionpoint{把屋顶模型的分母换成原子操作。}传统 roofline 用 FLOP/byte 描述每搬运一字节能做多少计算，本题则可定义 $I_a=5$ FLOP/atomic，原子屋顶为 $P_a=I_aR_a=0.05$ GFLOP/s。完整设备模型还应取 $P=\min(P_{\mathrm{peak}},I_bB,I_aR_a)$，分别表示计算峰值、普通内存带宽和原子服务率；题目只给出最后一项，因此答案是条件上界而不是 GPU 的总体能力。这个写法的价值在于把“很慢”变成可比较的资源预算，也可推广到每次锁获取、消息发送或稀疏索引查找所支撑的有效工作量。

\discussionpoint{指令能否重叠取决于数据依赖。}若五次浮点运算只使用早已准备好的寄存器，调度器可以在某个 warp 等待原子时发射其他 warp 的算术，因而单次原子延迟未必全部暴露在时间线上。反之，若运算读取原子返回值，关键路径会形成“原子、五次运算、下一原子”的串行链；此时增加独立线程只能在资源允许时交错多条链，不能消掉链内依赖。用依赖图而非简单指令计数看问题，可以解释为何两段拥有相同 FLOP 和原子次数的代码性能仍可能相差很大，也连接到指令级并行、软件流水和异步访存的基本方法。可用两个微基准分别让算术依赖或不依赖原子返回值，二者的差距就是关键路径作用的直接证据。

\discussionpoint{聚合会重新定义每次原子的有效工作量。}假设一个 warp 中有 $u$ 个线程更新同一 key，先用 ballot 和 shuffle 求和后只提交一次全局原子，理想情况下 $I_a$ 可从 5 增到 $5u$。但聚合本身需要投票、交换和分组，若 32 个线程的 key 全不相同，原子数没有减少，额外指令反而成为纯开销；线程块私有化也还要支付共享空间、同步和最终提交。因而收益由 key 的重复度和偏斜决定，这与图顶点度分布、稀疏梯度重复索引以及哈希聚合中的局部性直接相关，不能只用平均线程数推断。若聚合固定成本相当于四次普通操作，就至少需要足够重复 key 省下超过该成本的原子请求才可能获利。

\discussionpoint{有效 FLOP 才能回答业务生产率。}分析器可能统计地址计算、重试路径或编译器生成的浮点指令，但这些并不都代表算法向答案推进的五次运算。以图 frontier 为例，重复边造成的原子竞争会增加硬件工作，却不增加新发现顶点；以梯度累加为例，数值精度和确定性又可能比峰值操作数更重要。可靠报告应从算法定义给出有效 FLOP，再同时记录原子请求、普通字节流量和端到端时间，并用输入偏斜分层展示结果。这样的口径能避免把无效忙碌误写成高性能，也便于比较私有化、排序归约和直接原子三类完全不同的实现。
\variantheading 每原子改为 12 FLOP 时上限为 $1.2\times10^8$ FLOP/s，即 0.12 GFLOP/s。
\practiceheading 在 Nsight Compute Roofline 中把 atomic throughput 与 FLOP 计数联合核对。
\sourcecheck{https://docs.nvidia.com/cuda/cuda-c-programming-guide/}{NVIDIA CUDA C++ Programming Guide 的原子串行化语义；单位约分与十进制前缀独立复核，置信度高}
\end{solutionexercise}

\begin{solutionexercise}{4}
\problemheading 在第 1 题的条件下，假设内核把全局内存变量私有化为共享内存变量，共享
内存访问延迟为 1 ns。原来的全局内存原子操作全部改成共享内存原子操作。为简化
分析，假设把私有变量累加到全局变量所增加的全局内存原子操作，会使总执行时间
增加 10\%。又假设内核每执行一次原子操作就执行 5 次浮点运算。受原子操作吞吐量
限制时，内核执行能达到的最大浮点吞吐量是多少？

\approachheading 本章模型把一次原子操作近似为一次加载加一次存储，因此 1 ns 的
共享内存访问延迟对应 2 ns 原子服务时间；再把合并开销造成的总时间乘 1.1。

\answerheading 不计合并时，共享原子上限为
$1/(2\ \mathrm{ns})=5\times10^8$ ops/s，对应 $2.5\times10^9$ FLOP/s。加入
$10\%$ 时间开销后：
\[
P=\frac{2.5\times10^9}{1.10}
=2.2727\times10^9\ \mathrm{FLOP/s}
\approx2.27\ \mathrm{GFLOP/s}.
\]

\conclusionheading 最大约为 2.27 GFLOP/s。
\discussionheading
\discussionpoint{私有化是一棵分层归约树。}把全局计数器复制到每个线程块，相当于先让局部更新在近处汇合，再把每块结果送回根节点；原来每个输入一次的远端竞争因此变成每个非空分箱至多一次提交。还可继续设置每 warp 副本，甚至让线程先在寄存器中累计短游程，形成寄存器、warp、块、全局四级结构。层次越深，热点越容易被拆散，但初始化、同步、额外存储和合并轮次也越多；这与多核 CPU 的 per-core counter 和分布式 map--reduce 都体现相同原则，即通过增加中间状态换取更短的共享依赖链。

\discussionpoint{延迟模型不能代替吞吐模型。}题目把共享原子抽象成一次 1 ns 读加一次 1 ns 写，所以得到 2 ns；真实硬件可能由专门原子管线完成，并能让访问不同 bank 的请求流水重叠。于是“某次操作从发出到完成要多久”和“饱和后每秒能完成多少次”不再互为简单倒数，尤其当多个 warp 同时访问不同地址时更明显。相反，若所有 lane 命中同一地址或同一 bank，序列化会重新显露依赖链。理解这一区别有助于避免用指针追逐测得的延迟预测流式吞吐，也解释了缓存、Tensor Core 和网络交换机中普遍存在的 latency--throughput 分离。

\discussionpoint{用成本方程寻找私有化盈亏点。}设直接方案需要 $N$ 次全局原子，私有化后有 $N$ 次共享更新和 $S$ 次全局提交，则可比较 $T_d\approx NL_g$ 与 $T_p\approx NL_s+SL_g+T_{init}+T_{sync}$。题目把后几项合并成额外 10\% 时间，便于算出 2.27 GFLOP/s；现实中的 $S$ 取决于块数和每块非空分箱数，绝不是固定比例。低竞争、小输入或分箱极多时，清零与提交可能吞掉收益；高偏斜、大输入时，节省的全局序列化则通常占主导。类似的盈亏方程也可指导缓存分块、局部哈希表和批量消息合并。

\discussionpoint{共享容量会反过来限制并行度。}每块保存 $K$ 个 32 位计数需要 $4K$ 字节，若再为多个 warp 建副本，容量按副本数成倍增长；共享内存过大时，每个 SM 可驻留块数下降，隐藏延迟的能力也随之减弱。使用 16 位计数虽能节省空间，却必须证明单块命中数不会溢出，或设计定期冲刷；分批处理大直方图则会重复读取输入。生产实现因此要在均匀、Zipf 偏斜和单热点数据上同时观察局部冲突、全局提交与驻留量，并以端到端时间决定设计，而不能孤立追求最高共享原子吞吐。例如 $K=4096$ 的单副本已需 16 KiB，八个 warp 各建副本就达到 128 KiB，许多配置根本无法同时驻留。
\variantheading 若合并开销为 $25\%$，则吞吐为 $2.5/1.25=2.0$ GFLOP/s。
\practiceheading 用共享原子 microbenchmark 和 Nsight Compute 的 shared atomic/bank 指标测量目标 SM。
\sourcecheck{https://docs.nvidia.com/cuda/cuda-c-programming-guide/}{NVIDIA CUDA C++ Programming Guide 对共享/全局原子读改写的说明；对抗审查同时计入读写两次访问，并确认时间增加 10\% 应使吞吐除以 1.1}
\end{solutionexercise}

\begin{solutionexercise}{5}
\problemheading 某直方图内核处理包含 524,288 个元素的输入，生成包含 128 个分箱的直方图。
内核配置为每块 1024 个线程。
\begin{enumerate}[label=\alph*.,leftmargin=2.2em]
  \item 图 9.6 的内核不使用私有化、共享内存和线程粗化，它会在全局
  内存上执行多少次原子操作？
  \item 图 9.10 的内核使用私有化和共享内存，但不使用线程粗化。它在
  全局内存上最多可能执行多少次原子操作？
  \item 图 9.14 的内核使用私有化、共享内存和线程粗化，粗化因子为 4。
  它在全局内存上最多可能执行多少次原子操作？
\end{enumerate}
图 9.6 每个输入元素直接对全局分箱做一次原子更新；图 9.10 每块先在共享内存形成
128 个私有分箱，再把每个非零分箱至多原子累加到全局一次；图 9.14 的提交方式相同，
但每块由线程粗化处理 $4\times1024$ 个输入元素。

\approachheading 基本内核每输入一次全局原子；私有化内核每块最终对每个非零私有
分箱至多提交一次全局原子，因此最大值是块数乘 128。粗化会把块数缩小 4 倍。

\answerheading
\begin{enumerate}[label=\alph*.,leftmargin=2.2em]
\item 原书图 9.6 每个输入元素执行一次全局原子，共
$524{,}288$ 次。
\item 无粗化时块数为 $524{,}288/1024=512$。每个块的 128 个私有分箱都可能非零，
故全局提交至多 $512\times128=65{,}536$ 次。
\item 粗化因子 4 时，每块处理 $4\times1024=4096$ 个输入，块数为
$524{,}288/4096=128$；至多执行 $128\times128=16{,}384$ 次全局原子。
\end{enumerate}
私有直方图内部仍对输入执行共享内存原子；上面的 (b)、(c) 只统计题目所问的全局
内存原子。若某块没有命中某分箱，实际提交次数会低于上界。

\complexity{$O(524{,}288)$ 输入工作}{$128$ 个共享分箱/块}{块数从 512 降至 128；块内更新有分箱争用}
\conclusionheading (a) 524,288；(b) 最多 65,536；(c) 最多 16,384 次全局原子。
\discussionheading
\discussionpoint{上界和典型提交量回答不同问题。}每块 128 个分箱都非空时，私有化方案确实提交 $128$ 次全局原子，但随机输入未必达到这个上界。若块内 $m$ 个样本独立且均匀落入 $K$ 个箱，某箱为空的概率是 $(1-1/K)^m$，所以非空箱期望为 $K[1-(1-1/K)^m]$；例如 $m$ 很小时，它更接近 $m$ 而不是 $K$。这个占位模型还能用非均匀概率 $p_k$ 推广成 $\sum_k[1-(1-p_k)^m]$，从而把真实偏斜映射到提交量，连接到缓存工作集估计、布隆过滤器占用分析和数据库 distinct 计数。

\discussionpoint{私有化只是把竞争搬到更近的位置。}每个输入仍要更新一个共享分箱，因此单热点数据会让同一块的线程在共享地址上排队，只是远端全局竞争被隔离成多个局部队列。若 warp 中许多 lane 命中同一箱，可先用 ballot 找出同组 lane，再由一个代表执行原子；也可为每个 warp 建独立直方图，最后在块内归约。另一条路线是先按分箱排序再计算游程长度，它几乎消除原子却增加排序流量。三种方法分别适合高重复、容量充足和可复用有序结果的场景，说明优化必须匹配数据分布而非只匹配算子名称。可以先抽样估计每 warp 的唯一 key 数，再据此选择直接原子或聚合路径，避免为均匀数据支付无效分组。

\discussionpoint{粗化同时改变块数和块内覆盖率。}因子 4 让每块处理 4096 项，块数由 512 降到 128，这是全局提交上界缩小四倍的直接原因；但每块样本更多，也提高了其触及全部 128 个箱的概率。均匀模型下某箱仍为空的概率约为 $(127/128)^{4096}$，已经极小，因此本题的上界几乎就是期望值；若分箱数扩大到数万，结论便可能反转。这个双重效应与批处理大小选择相似：更大批次摊薄固定开销，却扩大局部工作集并减少并行任务数，需要把两条变化同时放进模型。在偏斜分布下，应把统一的 $1/128$ 换成各箱概率 $p_k$，稀有箱的空概率会远高于热门箱。

\discussionpoint{不同业务会把同一计数问题变成不同算法。}图像颜色直方图的箱数小且固定，共享数组通常合适；数据库 group-by 的键空间可能巨大，需要局部哈希表或排序；词频统计又常呈长尾分布，少数热词适合专门缓存、冷词则走普通路径。验证这些实现时，首先要用总和等于输入数等不变量保证没有丢计数，再观察每块非空箱分布和端到端吞吐，判断时间花在更新、初始化还是合并。这样从题目的三个确定计数出发，就能建立一个按键空间、偏斜度和结果复用需求选择算法的完整决策框架。若结果随后需要按 key 有序输出，排序归约的前期成本还能被下游直接复用，这会改变原本只看直方图阶段的选择。
\variantheading 若分箱数为 64，其余不变，则 (b) 最多 32768 次，(c) 最多 8192 次全局原子。
\practiceheading 用 Nsight Compute 统计 global/shared atomic transactions，并以均匀、单热点和真实输入基准测试。
\sourcecheck{https://docs.nvidia.com/cuda/cuda-c-programming-guide/}{NVIDIA CUDA C++ Programming Guide 的共享/全局原子作用域；按原书图 9.6、9.10、9.14 的提交循环逐块计数复核}
\end{solutionexercise}
